\documentclass[11pt]{article}
\usepackage[margin=.76in,top=.62in,bottom=.62in]{geometry}
\usepackage{amsmath,amssymb,amsthm,hyperref}
\renewcommand{\normalsize}{\fontsize{10}{12}\selectfont}
\DeclareMathSizes{10}{10}{8}{6}
\hypersetup{colorlinks=true,linkcolor=blue,urlcolor=blue,citecolor=blue}
\setlength{\emergencystretch}{2em}
\newtheorem{theorem}{Theorem}
\newtheorem{lemma}[theorem]{Lemma}
\newcommand{\R}{\mathbb R}
\newcommand{\Z}{\mathbb Z}
\begin{document}
\begin{center}
{\Large\bfseries An odd-total lower bound for simple SOC representations}\\[.5em]
Research proof and Lean formalization\\
October 26, 2026
\end{center}
\begin{abstract}
We answer the second open question in Section 7 of Wang's 2024 paper on
weighted geometric means: if the positive primitive integer weights have
odd sum, every simple second-order cone representation has at least as many
constraints as original input variables. The proof treats the literal
existential representation, allowing repeated indices, self-loops, unused
variables, and cycles initially. Equality in the previous counting bound
forces an injective input map. Logarithmic balance equations and a finite
doubling argument then force the total weight to be even. The full
representation theorem is formally verified in Lean.
\end{abstract}

\section{The source question}
Let $s_1,\ldots,s_m$ be positive integers and $S=\sum_j s_j$. Consider
\begin{equation}\label{gm}
 t^S\le\prod_{j=1}^m x_j^{s_j},\qquad x_j,t\ge0.
\end{equation}
Following Wang's Definition 3.1~\cite{Wang}, a simple SOC representation of
size $n\ge1$ has $n$ internal variables $y_1=t,y_2,\ldots,y_n$, the last
$n-1$ being existential nonnegative auxiliaries. For each $i$ it imposes
one inequality
\begin{equation}\label{constraints}
 y_i^2\le v_{c(i,0)}v_{c(i,1)},
\end{equation}
where $v$ lists all $m+n$ terminal and internal variables and the input map
$c$ is arbitrary. Its projection must be exactly~\eqref{gm}. Write
$L(s_1,\ldots,s_m)$ for the minimum size. Wang proved $L\ge m-1$
(Theorem 3.6) and asked on p.~1512 whether odd $S$, for primitive weights,
always gives $L\ge m$.

\begin{theorem}\label{main}
For positive integers $s_1,\ldots,s_m$ with odd sum, every simple SOC
representation of~\eqref{gm} has size $n\ge m$.
\end{theorem}
No gcd assumption is needed. In particular this proves the requested
statement when $\gcd(s_1,\ldots,s_m)=1$.

We use only soundness: any nonnegative assignment satisfying all
constraints must satisfy~\eqref{gm}. The already resolved bivariate
minimum-cardinality conjecture (Wang's Corollary 4.5) is a different
question and is not a new result here.

\section{Counting the input slots}
Draw an edge from each internal vertex to each of its two child slots in
\eqref{constraints}. Keep both slots even when their input indices coincide,
and let $r$ be the output vertex. Every terminal is reachable from $r$.
Indeed, otherwise assign value $1$ to all reachable vertices and $0$ to
all others. Each reachable internal vertex has both children reachable,
so its constraint holds with equality. Each unreachable internal vertex
has $y_i^2=0$ and a nonnegative product on the other side.
Thus every constraint holds, while $t=1$ and a positive-weight input is
zero, contradicting soundness.

If $R$ internal vertices are reachable, all $m$ terminals are reachable,
and each reachable vertex except $r$ is hit by a reachable child slot.
Consequently
\[
 m+R-1\le 2R,\qquad m-1\le R\le n.
\]
This reproves the old counting bound without irredundancy assumptions.
Suppose for contradiction that $n<m$. Necessarily $n=R=m-1$.
The map from all $2n$ child slots plus one extra slot for $r$ onto all
$m+n=2n+1$ vertices is now a bijection. Hence the child map is injective,
every nonroot vertex is used exactly once, and $r$ is never an input.

\section{Balance and the parity obstruction}
Let $A_{ij}$ count the child slots of internal vertex $i$ equal to internal
vertex $j$, and let $B_{ij}$ count its terminal children equal to terminal
$j$. Set $M=2I-A$. Every column of $A$ has sum at most $1$.
The matrix $M$ is invertible: if $M^T w=0$, then each $2w_j$ is either
zero or a single entry of $w$. Taking a largest absolute entry gives
$2\max_j|w_j|\le\max_j|w_j|$, so $w=0$.

For any $u\in\R^m$, solve $Mz=Bu$. The assignment of terminal values
$e^{u_j}$ and internal values $e^{z_i}$ satisfies every constraint with
equality. Soundness and strict monotonicity of the exponential give
\[
 S z_r\le\sum_j s_j u_j.
\]
The potential $(-u,-z)$ is also a solution, so equality holds for all $u$.
Define the real flow vector $q$ by
\[
 M^Tq=Se_r.
\]
Since $Mz=Bu$, we have $q^TBu=Sz_r=s^Tu$ for every $u$, whence
$q^TB=s^T$. The unique-input conclusions identify these equations as
\begin{align}
 q_i&=s_j &&\text{if a child of $i$ is terminal $j$},\label{leaf}\\
 q_i&=2q_j &&\text{if a child of $i$ is internal $j$},\label{internal}\\
 2q_r&=S.\label{root}
\end{align}
In~\eqref{internal}, $j\ne r$ because the root is not an input. Positivity
of the flow is neither assumed nor needed.

Let $D=\{i:q_i\notin\Z\}$. A child of a vertex in $D$ cannot be terminal,
by~\eqref{leaf}. If it is internal $j$, then~\eqref{internal} implies
$j\in D$, since an integer $q_j$ would make $q_i$ an integer. The injective
child map therefore restricts to an injection
\[
 D\times\{0,1\}\longrightarrow D.
\]
Finiteness yields $2|D|\le|D|$, so $D$ is empty. In particular $q_r$ is
an integer. Equation~\eqref{root} makes $S$ even, contrary to the
hypothesis. This proves Theorem~\ref{main}.

\section{Formal certificate and scope}
The certificate contains 41 declarations in seven modules, checked using
Lean 4.34.1 and mathlib revision \texttt{d13f23b7}. Its endpoint is
\begin{center}
\texttt{WGMParity.source\_odd\_total\_bound}.
\end{center}
The formal representation is the existential projection of the literal
nonnegative-real quadratic inequalities. Arbitrary finite index maps are
allowed; injectivity, root exclusion, matrix invertibility, potential
existence, and flow integrality are all proved. The stronger soundness-only
theorem is also certified. A fresh isolated replay uses warnings as errors,
has no proof holes or additional axioms, and reports only the standard
axioms \texttt{propext}, \texttt{Classical.choice}, and \texttt{Quot.sound}.
An independent exact enumeration of all 22,183 unordered input
configurations for $2\le m\le4$, $n=m-1$, including repeated inputs and
cycles initially, agrees with the theorem; this finite check is not used
in its proof.

The original $m-1$ bound and the general correspondence with mediated
sets are credited to Wang. The later graph-based optimization work
\cite{Blanco2024,Blanco2026} supplies broader modeling and computational
frameworks; our proof works directly with the existential inequalities.
The documented prior-work screen is bounded, and no absolute priority
claim is made. Independent review and final packet acceptance remain
pending for this draft.

\begin{thebibliography}{9}
\bibitem{Wang} J. Wang, \emph{Weighted geometric mean, minimum mediated set,
and optimal simple second-order cone representation}, SIAM J. Optim.
\textbf{34} (2024), 1490--1514.
\href{https://doi.org/10.1137/22M1531257}{doi:10.1137/22M1531257}.
\bibitem{Blanco2024} V. Blanco and M. Mart\'inez-Ant\'on,
\emph{On minimal extended representations of generalized power cones},
SIAM J. Optim. \textbf{34} (2024), 3088--3111.
\href{https://doi.org/10.1137/23M1617205}{doi:10.1137/23M1617205}.
\bibitem{Blanco2026} V. Blanco and M. Mart\'inez-Ant\'on,
\emph{On the structure and computation of extremal mediated graphs with
applications to conic optimization}, J. Optim. Theory Appl.
\textbf{210} (2026), article 24.
\href{https://doi.org/10.1007/s10957-026-03048-5}{doi:10.1007/s10957-026-03048-5}.
\end{thebibliography}
\end{document}
